%!TEX TS-program = xelatex
%!TEX encoding = UTF-8 Unicode

% Um dieses Dokument zu erzeugen sind folgende Pakete erforderlich:
% - texlive-xetex (oft aber nicht immer Teil der TexLive-Distribution)
% - texlive-lang-german (auch nicht immer dabei)
% - Möglicherweise mehr TexLive-Pakete (siehe Logdatei nach dem Kompilieren).
% - Die Schrift "Linux Libertine".
% - Das Paket "libertine.sty". Achtung, nicht mit dem gleichnamigen LaTeX-Paket
%		verwechseln, welches unglücklicherweise standardmässig auch von XeTeX
%		gezogen wird. Es genügt das mitgelieferte Paket im gleichen Ordner
%		wie das Tex-Dokument bereitzustellen.
% - Die Schrift "Liberation Mono"
% - Das FHNW-Logo "fhnw-technik-head.eps"
%
% Das Dokument wird erzeugt mit: $ xelatex Report.tex
% Aufgrund der Querverweise sind nach Änderungen jeweils zwei xelatex-Läufe nötig.
%
\documentclass[paper=A4,twoside=false,BCOR=0mm,DIV=calc,fontsize=12pt,enabledeprecatedfontcommands]{scrartcl}
%
\usepackage[automark,headsepline]{scrpage2}
\usepackage{xunicode,fontspec,xltxtra}
\usepackage[english,german]{babel}
\usepackage{graphicx}
\usepackage{lastpage}
\usepackage[pagebackref=true,colorlinks=true,allcolors=black]{hyperref}
\usepackage{listings}
\usepackage{csquotes}
%
\usepackage{enumitem}
\setlist[itemize]{noitemsep}
%
% Beautiful tables
\usepackage{booktabs}
\newcommand{\ra}[1]{\renewcommand{\arraystretch}{#1}} % Set on every table the space
%
% Graphs
\usepackage{pgf}
\usepackage{tikz}
\usetikzlibrary{arrows}
\usepackage{smartdiagram}
\usepackage{color}
\usepackage{xcolor}
%
\usepackage{pdfpages}
\usepackage{rotating}
%
\usepackage{pgfplots}
\usepgfplotslibrary{ternary}
\pgfplotsset{compat=1.8}
%
\usepackage{mathtools}
%\everymath{\displaystyle} use $\displaystyle \sum$ in text
\usepackage{amstext}
\renewcommand{\listfigurename}{Abbildungen}
%
% fuer Zitate
\usepackage{cite}
%\bibliographystyle{alphadin}
%\usepackage[debug]{libertine}
%\setromanfont[Mapping=tex-text]{Linux Libertine O}
%\setsansfont[Mapping=tex-text]{Linux Biolinum O}
%\setmonofont[Mapping=tex-text]{Liberation Mono}

\pagestyle{scrheadings}
\clearscrheadfoot{}
\ihead{\headmark}
%\ohead{Seite\pagemark\ von \pageref{LastPage}}
%\ifoot{P. Steinger \& N. Mauchle}
%\ofoot{\today}
\ofoot{Seite\pagemark\ von \pageref{LastPage}}

% Two images beside
\usepackage{subcaption}
%
%
\setlength{\parindent}{0pt}
\setlength{\parskip}{1em}
%
\newenvironment{myenumerate}{
\begin{enumerate}
  \setlength{\itemsep}{0pt}
  \setlength{\parskip}{0pt}
}{\end{enumerate}}
%
\newenvironment{myitemize}{
\begin{itemize}
  \setlength{\itemsep}{0pt}
  \setlength{\parskip}{0pt}
}{\end{myitemize}}
%
\renewcommand*{\pnumfont}{
	\normalfont\rmfamily\slshape
}
%
\KOMAoptions{draft=true}
\KOMAoptions{DIV=last}
%
\textheight=630pt
%
%\usepackage{layout}
\begin{document}
\overfullrule=0pt
%
\definecolor{lightgray}{rgb}{.9,.9,.9}
\definecolor{darkgray}{rgb}{.4,.4,.4}
\definecolor{purple}{rgb}{0.65, 0.12, 0.82}
%
\lstdefinelanguage{Python}{
  keywords={typeof, new, True, False, except, function, return, null, catch, switch, var, if, in, while, do, else, case, break},
  keywordstyle=\color{blue}\bfseries,
  ndkeywords={class, export, boolean, throw, implements, import, this},
  ndkeywordstyle=\color{darkgray}\bfseries,
  identifierstyle=\color{black},
  sensitive=false,
  comment=[l]{//},
  morecomment=[s]{/*}{*/},
  commentstyle=\color{purple}\ttfamily,
  stringstyle=\color{red}\ttfamily,
  morestring=[b]',
  morestring=[b]"
}
%
\lstset{
   language=Python,
   backgroundcolor=\color{lightgray},
   extendedchars=true,
   basicstyle=\footnotesize\ttfamily,
   showstringspaces=false,
   showspaces=false,
   numbers=left,
   numberstyle=\footnotesize,
   numbersep=9pt,
   tabsize=2,
   breaklines=true,
   showtabs=false,
   captionpos=b
}
%
%\layout
%
% --- Titelseite --- %
\begin{titlepage}
	\enlargethispage{3cm}
	\begin{raggedright}
	\begin{picture}(0,0)
		\put(-30,14){\includegraphics[width=7cm]{fhnw-technik-head}}
	\end{picture}

	\vspace*{4cm}
	{\Huge\bfseries\normalfont\sffamily
		Mit Machine Learning Immobilienpreise schätzen\\[1.7ex]
	}
	{\Large\bfseries\normalfont\sffamily
		Vergleich von verschiedenen Machine Learning Algorithmen um Immobilienpreise zu schätzen.\\[2.2ex]
	}
	{\large\bfseries\normalfont\sffamily
		Projekt 6 von\\[1.5ex]
		Piero Steinger\\[1.5ex]
		Nicolas Mauchle\\[1.5ex]
	}
	\vspace*{1.5cm}
	{\large\bfseries\normalfont\sffamily
		FHNW\\[1.5ex]
		Hochschule für Technik\\[1.5ex]
		Studiengang IT\\[2.5ex]

		Examinator:\\[1.5ex]
		Prof. Dr. Manfred Vogel\\[2.5ex]

		Experte:\\[1.5ex]
		Johanes Schwertfeger\\[1.5ex]
	}
	\vspace*{2cm}
	{\large\bfseries\normalfont\sffamily
		Windisch, \today\\
	}
	\end{raggedright}
\end{titlepage}
%
\newpage
%
\pagenumbering{Roman}
% --- Zusammenfassung --- %
\section*{Abstract}
Häuserpreise werden Heute von Experten geschätzt. Dafür muss ein Experte das Haus begutachten und diverse Kennwerte aufnehmen. Das kostet Geld und ist für beide Parteien mit Aufwand verbunden. Zudem weichen Schätzungen von verschiedenen Makler zum Teil stark voneinander ab.\\
In diesem Projekt wird analysiert, ob ein Machine Learning Ansatz nicht ähnliche oder sogar bessere Schätzungen abgeben kann. Zu diesem Zweck werden Immobiliendaten gesammelt und analysiert. Anhand von verschiedenen Machine Learning Algorithmen wird versucht ein exaktes Schätzungsmodell zu berechnen.
% --- Vorwort ----------------------------------------------------- %
\section*{Vorwort}
Das Projekt \emph{Mit Machine Learning Immobilienpreise schätzen} ist eine Forschungsarbeit im Rahmen der Bachelor Arbeit. Es soll untersucht werden, ob öffentliche Immobiliendaten ausreichen um akkurate Schätzungsmodelle zu erstellen.\\[2ex]
Die Autoren bedanken sich bei Prof. Dr. Vogel, er hat das Projekt betreut und unterstützt. 


%
% --- Inhaltsverzeichnis --- %
\newpage
	\tableofcontents
%
\clearpage
\pagenumbering{arabic}
\newpage
%
% --- Einleitung --- %
\section{Problemstellung}
Machine Learning wird heute in sehr vielen Gebieten eingesetzt \cite{forbes}. Das Schätzen eines Kaufpreises einer Immobilie ist ein klassischer Anwendungsfall. Der Grund dafür sind die vielen Kennwerte, die eine Immobilie besitzt.\\
Heute werden Immobilienpreise\footnote{Der Verkehrswert einer Immobilie.} immer noch von Experten geschätzt. Dies obwohl Machine Learning Algorithmen in der Lage sind eine gleichwertige oder sogar präzisere  Schätzung abzugeben.\\
Damit ein gutes Schätzungsmodell berechnet werden kann, braucht es einen aussagekräftigen Datensatz. Einen solchen Datensatz mit genügend Daten zu erstellen ist nicht trivial und kostet viel Zeit. Ausserdem existiert noch kein öffentlich zugänglicher Datensatz für den schweizer Immobilienmarkt.\\
Für die Beschaffung der Immobiliendaten in der Schweiz bieten sich daher die vielen Immobilienplattformen im Internet an. Leider sind diese Daten stellenweise unvollständig oder fehlerhaft erfasst. Auch können Immobilien über zwei oder mehrere Plattformen beworben werden.\\
Es ist daher interessant zu wissen, ob mit öffentlichen Daten einen Datensatz erstellt werden kann, der mit aktuellen Machine Learning Algorithmen den Immobilienpreis akkurat schätzt. Zudem wäre dies auch für Makler und Immobilienbesitzer attraktiv, da sie diese Schätzungen als Unterstützung verwenden können.
%
%
\subsection{Ausgangslage}
In der Schweiz konkurrieren dutzende Immobilienplattformen im Internet, auf denen Immobilien zum Kauf oder Miete ausgeschrieben sind. Täglich kommen neue Inserate hinzu. Diese Daten stehen jedem frei zur Verfügung. Sie sind jedoch von Plattform zu Plattform unterschiedlich erfasst und müssen für die weitere Verwendung harmonisiert werden. Darüber hinaus müssen fehlerhafte oder doppelte Inserate erkannt und entfernt werden. Ferner können ortsbezogene Informationen vom Bundesamt für Statistik (BFS) eingeholt werden, um ein Inserat mit zusätzlichen Informationen zu ergänzen. Die Daten vom BFS stehen ebenso öffentlich zur Verfügung.

Zum Thema “Kaufpreise für Immobilien schätzen” wurden diverse Forschungsarbeiten veröffentlicht (vgl. \cite{existing_work_1, existing_work_2, existing_work_4}).
Nicht immer ist der einzelne Kaufpreis einer Immobilie gefragt, sondern mehr die Entwicklung des Immobilienpreises für eine spezifische Ortschaft \cite{existing_work_5}.\\
Für die Schweiz gibt es ähnliche Forschungsarbeiten. Meist wird jedoch versucht den Mietpreis pro Quadratmeter und nicht den Kaufpreis zu schätzen \cite{existing_work_3, existing_work_6}.
%
\subsection{Hypothese und Forschungsfragen}
In dieser Arbeit soll erforscht werden ob mit öffentlich zugänglichen Immobilieninseraten und Machine Learning Algorithmen der Kaufpreis einer Immobilie akkurat geschätzt werden kann.\\
Zusätzlich zu den Immobilieninseraten werden geobasierte Daten vom Bundesamt für Statistik hinzugezogen um zu überprüfen ob eine präzisere Schätzung erzielt werden kann.

Um den Kaufpreis einer Immobilie zu bestimmen, können verschiedene Verfahren angewendet werden. Diese Verfahren entsprechen keiner exakten Wissenschaft. Unter anderem auch deshalb, weil der Kaufpreis stark von Angebot und Nachfrage abhängig ist. Experten auf dem Gebiet sprechen von einem wirtschaftlich verkraftbaren Preis. Somit unterliegen die Preise einer gewissen Schwankung. Bei einer Experteneinschätzung wird von einer übereinstimmenden Bewertung gesprochen, wenn die Schätzwerte nicht mehr als 10\% voneinander abweichen vgl. \cite{immo_1, immo_2}.\\[2ex]
%
Für die Arbeit wird folgende These aufgestellt:\\[2ex]
\textit{Anhand von öffentlich gesammelten Daten können drei Viertel aller Immobilienpreise mithilfe von Machine Learning Algorithmen mit einer maximalen Abweichung von 10\% geschätzt werden.}\\[4ex]
%
Um diese These zu beantworten, werden folgende Forschungsfragen untersucht:
\begin{enumerate}
\item Wie gut reichen öffentliche Immobiliendaten aus, um aussagekräftige Daten zu erhalten?
\item Wie stark kann das Resultat verbessert werden, wenn geobasierte Daten hinzugefügt werden?
\item Welche Machine Learning Algorithmen eignen sich am Besten, um Immobilienpreise zu schätzen?
\end{enumerate}
Als Endergebnis soll eine Softwareapplikation aufzeigen, wie genau der Immobilienpreis geschätzt werden kann. Dazu werden diverse Algorithmen mit verschiedenen Parametern untersucht. Dabei kann nicht ausgeschlossen werden, dass noch exaktere Varianten oder Algorithmen vorhanden sind, die ein genaueres Resultat erzielen.
%
\subsection{Methode}
Für diese Arbeit wird mit dem iterativen Vorgehensmodell Plan Do Check Act (PDCA) gearbeitet. Dies erlaubt nach jeder Iteration das Resultat zu analysieren und, wenn möglich, in einer weiteren Iteration zu verbessern. Dazu ist es flexibel genug um auf Fehlentscheide schnell reagieren zu können.
%
\subsection{Struktur}
Im nächsten Kapitel werden die verwendeten Machine Learning Algorithmen vorgestellt. Darauffolgend wird die Architektur aufgezeigt. Im Kapitel Validierung werden die gesammelten Daten ausgewertet und die Machine Learning Algorithmen untersucht um die Studienfrage zu beantworten. Im letzten Kapitel findet sich ein Résumé der Arbeit.
\newpage
%
\section{Hintergrund}
In diesem Kapitel wird auf die Konzepte und Begriffe von Webcrawlern und Machine Learning eingegangen. Es soll ein gemeinsames Grundwissen für beide Gebiete geschaffen werden.
%
\subsection{Webcrawler}
Um Informationen von Webseiten zu sammeln, muss deren Inhalt analysiert und die gewünschten Daten extrahiert werden. Für diese Aufgabe eignen sich Webcrawler besonders gut \cite{scrapy_1, scrapy_2}.\\
Dabei erhält ein Webcrawler eine oder mehrere Einstiegs-URLs. Diese werden dann vom Crawler aufgerufen und die dazugehörigen Seiten ausgewertet. Bei der Auswertung kann der Crawler auf weitere URLs stossen, die er in seine abzuarbeitende URL Liste aufnehmen kann. Sind alle URLs abgearbeitet ist der Crawler fertig. Je nach Aufgabenbereich existieren verschiedene Crawler-Techniken.\\
Ist bekannt von welchen Seiten die Informationen gecrawlt werden können, eignet sich am besten ein fokussierter Crawler. Dieser Crawler steuert vordefinierte Seiten an und extrahiert Informationen nach vorgegebenen Regeln.\\
Sind die Webseiten nicht bekannt oder kommt man nicht ohne weiteres an gesuchte Informationen, eignen sich Deep Web Crawler. Diese bewegen sich mehrheitlich autonom und versuchen durch Analysieren des Inhaltes herauszufinden ob der Inhalt von Interesse ist oder nicht.\\
Sind es übermässig viele Seiten, die gecrawlt werden müssen, oder ist die Stabilität der Crawler wichtig, kommen verteilte Crawler zum Einsatz. Durch die Verteilung wird an Stabilität gewonnen und die Last auf die verschiedenen Crawler verteilt \cite{scrapy_3}.\\[2ex]
%
Um Daten von einer Webseite zu extrahieren, gibt es diverse Möglichkeiten. Von der einfachen Stringsuche bis zu Machine Learning Methoden gibt es nahezu alles. Beim fokussierten Crawlen wird häufig versucht Informationen über den XPATH oder über die CSS-Klassen zu extrahieren. Beide Varianten haben ihre Vor- und Nachteile, sind aber im Grunde sehr ähnlich \cite{xpath_vs_css}. Somit ist es dem Benutzer überlassen welche Variante er wählt \cite{xpath}.\\[2ex]
%
Eine Herausforderung beim Crawlen ist, dass der Crawler nicht als solchen vom Webseitenbetreiber erkannt wird. Diverse Betreiber möchten gerne verhindern, dass Informationen von Ihren Seiten gecrawlt werden \cite{comparis}.\\
Ist ein Crawler zu aggressiv, kann es vorkommen, dass er vom Betreiber gesperrt wird. Es gilt daher mit einer gewissen Raffinesse vorzugehen. Um möglichst nicht aufzufallen, ist es sinnvoll nur ein paar Requests\footnote{Anfrage auf eine Webseite} pro Minute abzusetzen und einen existierenden User Agent anzugeben. Wenn das nicht ausreichend ist, können die Seiten über einen Pool von rotierenden IP-Adressen aufgerufen werden. So ist die Quelladresse nicht immer dieselbe \cite{offensive_crawling}.
%
\subsection{Machine Learning}
Machine Learning ist eine Art Künstliche Intelligenz (AI), das Softwareapplikationen erlaubt Vorhersagen zu machen, anhand von trainierten Algorithmen.
Die Definition von Machine Learning aus dem Jahr 1959 von Arthur Samuel, einer der Pioniere im Bereich Machine Learning, lautet \cite{what_is_ml}: 
  \begin{quote}
    Machine Learning: Field of study that gives computers the ability to learn without being explicitly programmed.
  \end{quote}
Mit Hilfe von Machine Learning Algorithmen können neben Vorhersagen auch Trends erkannt oder Klassifizierungen durchgeführt werden \cite{ml_book}.
Dabei handelt es sich nicht um eine exakte Wissenschaft. Es lohnt sich verschiedene Algorithmen mit unterschiedlichen Parametern auszuprobieren, um die geeignetste Performance zu finden. Vieles hängt auch von den Ausgangsdaten ab \cite{ml_azure}.\\
%
Grundsätzlich erhält ein Algorithmus Daten, die er zu einem Modell verarbeitet. Dieses Modell wird ausgewertet. Entspricht es nicht den gewünschten Anforderungen wird versucht ein besseres Modell zu erstellen, bis das Resultat zufriedenstellend ist. Abbildung \ref{fig:ml_process} zeigt den schematischen Ablauf.\\
\begin{figure}[ht]
\centering
\begin{tikzpicture}[
  font=\sffamily,
  every matrix/.style={ampersand replacement=\&,column sep=1.5cm,row sep=.5cm},
  source/.style={draw,thick,rounded corners,fill=yellow!20,inner sep=.3cm},
  process2/.style={draw,thick,rounded corners,fill=blue!20,inner sep=.3cm},
  process/.style={draw,thick,circle,fill=orange!20},
  process3/.style={draw,thick,rounded corners,fill=orange!20, inner sep=.3cm},
  sink/.style={source,fill=green!20},
  datastore/.style={draw,very thick,shape=datastore,inner sep=.3cm},
  dots/.style={gray,scale=2},
  to/.style={->,>=stealth',shorten >=1pt,semithick,font=\sffamily\footnotesize},
  dotted/.style={dashed,->,>=stealth',shorten >=1pt,semithick,font=\sffamily\footnotesize},
  every node/.style={align=center}]
% Position the nodes using a matrix layout
\matrix{
  \node[source] (data) {Data};
      \& \node[process] (model) {Model};
        \& \node[sink] (prediction) {Prediction}; \\
};

% Draw the arrows between the nodes and label them.
\draw[to] (data) -- node[midway,above] {Train}
    node[midway,below] {} (model);

\draw[to] (model) -- node[midway,above] {Predict}
    node[midway,below] {} (prediction);

\draw[to] (prediction) to[bend right=50] node[midway,above] {Evaluate}
    node[midway,below] {} (data);
\end{tikzpicture}
\caption[Schematischer Machine Learning Ablauf]{Schematischer Machine Learning Ablauf}%
\label{fig:ml_process}
\end{figure}
\newline
%
Es gibt drei verschiedene Kategorien von Machine Learning \cite{super_unsuper}.\\
\textbf{Supervised Learning:}	Bei diesen Algorithmen sind die Daten gekennzeichnet. Das heisst, es ist erkenntlich, welchen Beobachtungswert die Daten darstellen. Mit Hilfe von diesen Daten, wird ein Modell erstellt, das möglichst nahe den Beobachtungswert schätzt. Dabei kann es sich um eine Klassifizierung oder Regression handeln.\\[2ex]
%
\textbf{Unsupervised Learning:} Hier hat man keine gekennzeichneten Daten. Es wird versucht Muster im Datensatz zu erkennen. Anhand von diesen können unter anderem Trends oder Ausreisser erkannt werden.\\[2ex]
%
\textbf{Reinforcement Learning:} Hier lernt der Algorithmus anhand seiner Aktionen. Für jede Aktion, die er berechnet, wird er entweder belohnt oder bestraft. Mit der Zeit kann der Algorithmus situationsbedingte Abläufe erlernen und später anwenden.\\[2ex]
%
Für eine Schätzung braucht jeder Machine Learning Algorithmus ein berechnetes Modell. Dieses kann unter anderem eine mathematische Formel oder eine Baumstruktur sein. Damit ein Modell erstellt werden kann, braucht jeder Algorithmus ein oder mehrere Features. Features sind messbare Merkmale, die einen Beobachtungswert, auch Target genannt, beschreiben. Es ist sinnvoll die Features im Voraus zu untersuchen und herauszufinden, welche Features wichtig sind und welche weggelassen werden können. Dabei können zwei oder mehrere Features kombiniert werden. Dieses Vorgehen wird auch Feature Engineering genannt.\\[2ex]
%
Damit ein Algorithmus überprüfen kann, wie gut seine Schätzungen sind, wird der Fehler, beziehungsweise die Residuals berechnet. Die Differenz zwischen einem geschätztem Wert und dem dazugehörigen beobachteten Wert wird als Residual bezeichnet. Dabei sind systematische\footnote{Ein falsches Modell wurde verwendet.} wie auch statistische Fehler\footnote{Es handelt sich um Ausreisser.} im Residual mit  eingerechnet. Formel \eqref{eq:residuals} zeigt die Berechnung des Residuals, dabei ist $\hat{y}$ der geschätzte Wert und y der beobachtete Wert. Es kann sich dabei sowohl um Skalare als auch Vektorenwerte handeln.\\
\begin{equation}
\label{eq:residuals}
\hat{\varepsilon} = y - \hat{y}
\end{equation}
%
\begin{table*}[hb]
\centering
\ra{1.3}
\resizebox{\textwidth}{!}{
\begin{tabular}{@{}llll@{}}
\toprule
Name & Abkürzung & Formel & Modelart\\
\midrule
Mean Absolute Error & MAE & $\frac1n \sum_{i=1}^{n} |y_i - \hat{y_i}|$ & absolut\\
Median Absolute Error & MdAE & $med(\sum_{i=1}^{n} |y_i - \hat{y_i}|)$ & absolut\\
Sum of Squared Errors & SSE & $\sum_{i=1}^{n} (y_i - \hat{y_i})^2$ & quadratisch\\
Root Mean Squared Errors & RMSE & $\sqrt{\frac1n \sum_{i=1}^{n} (y_i - \hat{y_i})^2}$ & quadratisch\\
Mean Absolute Percentage Error & MAPE & $\frac1n \sum_{i=1}^{n} |\frac{y_i - \hat{y_i}}{y_i}|, y_1 \neq 0$ & prozentual\\
Median Absolute Percentage Error & MdAPE & $med(\sum_{i=1}^{n} |\frac{y_i - \hat{y_i}}{y_i}|), y_1 \neq 0$ & prozentual\\
\bottomrule
\end{tabular}}
\caption{Fehlermodelle mit Formelen}
\label{tab:error_models}
\end{table*}
%
Wie in Tabelle \ref{tab:error_models} gezeigt, haben sich mit der Zeit diverse Fehlermodelle entwickelt. Die Wahl des Fehlermodells ist vom Anwendungsfall abhängig. Grundsätzlich können Fehler absolut oder quadratisch berechnet werden \cite{error_models}.\\
Bei absoluten Werten werden alle Fehler gleich stark gewichtet, wobei bei quadratischen Fehler grössere Fehler durch das Quadrieren stärker gewichtet werden. Beim MAE hat ein Fehler bei einem hohen Beobachtungswert einen grösseren Einfluss auf den gesamten Fehler. Das kann durch den Median (MdAE) oder einer prozentualen Angabe wie der MAPE oder MdAPE vermieden werden. Eine prozentuale Angabe hat zusätzlich den Vorteil, dass sie intuitiver Interpretiert werden kann. Der Nachteil bei MAPE und MdAPE ist, dass sie bei einem y Wert von 0 nicht verwendet werden können \cite{error_models_2}.\\[2ex]
%
Eine weitere Möglichkeit die Performance eines Modells neben dem Fehlermodell zu bestimmen, ist der $R^2$ Wert. $R^2$ ist ein Gütemass, das beschreibt, wie gut die Features in der Lage sind, die Varianz der Zielvariable erklären zu können. Die Formel \eqref{eq:r2} zeigt, dass der Wert immer zwischen 0 bis 1 liegt. Dabei ist y der Beobachtungswert, $\hat{y}$ der geschätzte Wert und $\overline{y}$ der Durchschnitt aller y. Sind alle Beobachtungswerte gleich ist $R^2$ nicht definiert. Je näher sicher der Wert bei 1 befindet, desto besser ist das Modell.\\
Es gilt jedoch zu beachten, dass der $R^2$ Wert die Varianzen aufzeigt. Somit kann nicht davon ausgegangen werden, dass ein hoher $R^2$ Wert auch eine hohe Genauigkeit beim Schätzen erzielt \cite{r2, r2_2}. 
\begin{equation}
\label{eq:r2}
R^2 = 1 - \frac{\sum_{i=1}^{n} (y_i - \hat{y_i})^2}{\sum_{i=1}^{n}(y_i - \overline{y})^2},\text{ falls nicht }y_1  = ... = y_n = \overline{y}
\end{equation}
%
Es macht keinen Sinn auf denselben Daten zu trainieren und das Modell zu überprüfen. Deshalb wird der Datensatz in mindestens einen Trainings- und einen Testdatensatz aufgeteilt \cite{cross_validation}.
Dadurch kann die Gefahr eines Overfit verringert werden. Bei einem Overfit wurde das Modell zu stark an den Testdaten trainiert. Wichtig dabei ist, dass die Daten vor dem Aufteilen gründlich gemischt werden. So, dass nicht alle Wohnungen im Trainingsdatensatz sind und nicht alle Häuser im Testdatensatz.\\
Cross Validation ist eine weitere Möglichkeit den Datensatz aufzuteilen. Dieser teilt den gesamten Datensatz in K-Teile auf. Dabei werden K-1 Datensätze für das Trainieren und 1 Datensatz für das Testen verwendet. Das ganze wird K-mal wiederholt, so, dass jeder Datensatz einmal als Testdatensatz verwendet wurde.
%
\subsubsection{Lineare Regression}
Bei der Linearen Regression versucht man anhand einer linearen Funktion einen Beobachtungswert y durch einen oder mehrere Features X zu berechnen. Dabei handelt es sich um ein statistisches Verfahren bei dem davon ausgegangen wird, dass y abhängig von X ist. Weiter gilt, dass die unabhängigen Variablen X dichotom oder intervallskaliert sind.
\begin{equation}
\label{eq:linear}
\hat{y} = h_\theta(x) = \sum_{j=0}^{N} \theta_j x_j
\end{equation}
Formel \eqref{eq:linear} zeigt die Lineare Regression wobei  den Koeffizienten für jedes einzelne Feature darstellt. Je höher der Koeffizient ist, desto mehr Gewicht hat dieses Feature im Modell. Um auf die optimalen Koeffizienten zu kommen, versucht der Algorithmus eine Gerade durch den Datensatz zu ziehen, die den kleinsten Abstand zu allen Punkten hat. Sprich, der Algorithmus versucht das Schätzungsmodell auf das gewählte Fehlermodell zu optimieren. Zugleich dient das Fehlermodell auch als Performancemetrik. Die Optimierung wird durch die Kostenfunktion \eqref{eq:cost_function} berechnet. Oft wird dafür die \textit{Ordinary least squares} (OLS) Methode genommen. Es gibt auch noch weitere Methoden für die Optimierung.\\
Das quadratische Fehlermodell macht den Algorithmus anfälliger auf Ausreisser. Diese sollten, wenn immer möglich, in einem vorherigen Schritt ausgemustert werden. Auch sollte der Fehler in den Daten normalverteilt sein und die Residuen eine Homoskedastizität aufweisen, da ansonsten keine gute Regression gefunden wird \cite{gradient_descent, gradient_descent_2, gradient_descent_3}.
\begin{equation}
\label{eq:cost_function}
J(\theta) = \frac{1}{2m} \sum_{i=1}^{N} (h_\theta(x^i) - y^i)^2
\end{equation}
%
\textbf{Polynomial Regression}\\
In vielen Fällen hat der Datensatz keine Homoskedastizität, sondern besitzt eine Heteroskedastizität \cite{poly}. Dies ist oft dann der Fall wenn das Verhältnis zwischen y und X einer Kurve entspricht. In solch einem Fall, erzielt die Lineare Regression keine gute Performance.
Bei der Polynomialen Regression ist deshalb der beste Fit keine Gerade, sondern ähnelt einer Kurve. Je nachdem wie viele Polynome die Funktion besitzt, desto kurvenreicher ist die Funktion am Schluss.\\
Die Formel für eine Polynom Gleichung zweiten Grades wird in Gleichung \eqref{eq:poly} dargestellt.
%
\begin{equation}
\label{eq:poly}
y = \theta_1 + \theta_2 * X + \theta_3 * X^2
\end{equation}
%
\newline
Je höher der Grad des Polynoms, desto genauer wird der Algorithmus auf den Trainingsdatensatz trainiert. Das hat den Nachteil, dass neue Datensätze ungenau geschätzt werden, da das Modell eine hohe Varianz aufweist. Dieses Phänomen nennt sich Overfitting. Der Algorithmus wurde zu stark am Trainingsdatensatz trainiert.\\
Ist der Grad des Polynoms zu klein gewählt, kommt es zu einem Underfitting. Das Modell kann nur Features abbilden, die dem ausgewählten Grad entsprechen. Alle anderen Features werden ignoriert. So entsteht die Gefahr, dass wichtige Features nicht ins Modell miteinfliessen und das Modell einen hohen Bias besitzt.\\
Es wird auch von dem \textit{Bias-Varianz-Dilemma} gesprochen \cite{bias_variance, bias_variance_2}. Denn die Kunst ist, die richtige Abstimmung zu finden, damit das Modell die beste Performance hat.

\newpage
\clearpage
% NUR ZIEL SYSTEM (Das was wir entwickeln) WIE wird mit situationen gearbeitet, was wir machen
\section{Architektur}

\newpage
%
\clearpage
\section{Validierung}
Um die Forschungsfragen zu beantworten, haben wir die gesammelten Immobiliendaten auf Verwertbarkeit analysiert und gefiltert. Die bereinigten Inserate wurden für den Vergleich der verschiedenen Machine Learning Algorithmen verwendet. Dabei diente der MAPE sowie der MdAPE als wichtigste Performanzkriterien. Weiter war uns wichtig, wie viele Inserate mit einer maximalen Abweichung von 10\% geschätzt werden können. Mit Hilfe von Feature Engineering und Parameter Optimierung wurde versucht die Lernalgorithmen zu verbessern.\\
751 Inserate haben so extreme Werte, dass sie die Statistik verzerrten und so eine Aussage über die Daten erschwerten. Deshalb wurden diese Werte für die Validation ignoriert.

\subsection{Gesammelte Daten}
Insgesamt konnten 162’225 Immobilieninserate in einem Zeitraum von 5 Monaten gesammelt werden. Davon stammten die meisten vom Immoscout24 Portal, wie in Tabelle \ref{tab:crawled_data} ersichtlich.

\begin{table*}[ht]
\centering
\ra{1.3}
\begin{tabular}{@{}lrr@{}}
\toprule
Portal &  Anz. gesammelte Inserate & Anz. verwendete Inserate \\
\midrule
Immoscout24 & 94'611 & 46'226\\
Homegate & 34'115 & 19'568\\
Newhome & 32'600 & 19'011\\
Urbanhome & 2'047 & 1'040\\
\bottomrule
\end{tabular}
\caption{Auswertung der gesmmelten Daten}
\label{tab:crawled_data}
\end{table*}

Von den insgesamt 162’225 Inseraten wurden 77’608 ungültige Inserate herausgefiltert.
Dabei handelte es sich mehrheitlich um fehlerhafte Inserate oder Inserate, bei denen wichtige Kennwerte fehlten.\\
Insgesamt konnten 83’107 Wohungen und 65’430 Häuser gesammelt werden. Die restlichen Inserate waren Lagerräume oder nicht definierte Immobilien\-typen.\\
Weiter konnte für 99\% aller Inserate die Gemeinde zugeordnet, sowie für 83\% aller Inserate, mit einer vollständigen Adresse, eine Lärmbelastung berechnet werden.\\
Die meisten Inserate hatten eine Angabe über die Anzahl Zimmer und Wohnfläche, gefolgt vom Baujahr und der Anzahl an Stockwerken, falls es mehrere besass. Die Raumhöhe, die Kubatur und die Gesamtfläche einer Immobilie wurden dagegen nur selten angegeben. Abbildung \ref{fig:features} zeigt den prozentualen Anteil der fehlenden Kennwerte.

\begin{figure}[h!]
\centering
\includegraphics[width=0.9\textwidth]{images/missing_values.png}
\caption[\% Anteil der fehlenden Features]{\% Anteil der fehlenden Features}%
\label{fig:features}
\end{figure}

Die Crawler wurden nie über eine längere Zeit gesperrt. Vereinzelt kam es zu Blockierungen, die jedoch, durch das Starten neuer Proxyinstanzen, umgangen werden konnten. Der Zeitaufwand um ein Portal bei fünf gleichzeitigen Verbindungen und einer Downloaddelay von zwei Sekunden zu crawlen, lag im Schnitt bei 4 Stunden. Dementsprechend wurden in dieser Zeit etwa 100 Proxyinstanzen gestartet.\\[2ex]
%
Die Datenqualität der gesammelten Inserate war eher mässig, viele Features konnten nicht für das Machine Learning verwendet werden. Um dem entgegenzuwirken, wurde versucht durch Feature Engineering neue Features zu generieren.

\subsubsection{Kaufpreis}
Abbildung \ref{fig:price} zeigt die Verteilung der Kaufpreise. Abbildung \ref{fig:verteilung_price_all} beinhaltet alle Inserate wobei Abbildung \ref{fig:verteilung_price_part} nur die Inserate von 0 CHF - 3,5 Mio. CHF aufzeigt. Dabei ist ersichtlich, dass der Preis nicht normalverteilt ist. Es besteht eine positive Skenwess, was für eine lineare Regression nicht optimal ist.\\
Die Mehrheit der Preise befindet sich zwischen 250'000 CHF und 1.2 Mio. CHF. Daneben kommen mehrheitlich teurere Immobilien vor. Tabelle \ref{tab:price} zeigt die beschreibenden statistischen Werte des Kaufpreises vor dem Feature Engineering. Es ist zu sehen, dass die Standardabweichung von Wohnungen fast halb so gross ist, wie die der Häuser. Das liegt unter anderem daran, dass unsere gesammelten Wohnungspreise kein grosses Preisspektrum besitzen. Zum Vergleich werden im Kapitel Feature Engineering die beschreibenden statistischen Werte des Preises nach der Outlier Detection angezeigt.

\begin{figure}[ht]
\begin{subfigure}{.5\textwidth}
  \centering
  \includegraphics[width=\linewidth]{images/bar_des_kauf_preises.png}
  \caption[Verteilung des Preises mit allen Daten]{Verteilung des Preises mit allen Daten\\}
  \label{fig:verteilung_price_all}
\end{subfigure}
\begin{subfigure}{.5\textwidth}
  \centering
  \includegraphics[width=\linewidth]{images/bar_des_kauf_preises_cut.png}
  \caption[Verteilung des Preises von 0 CHF - 2.5 Mio CHF]{Verteilung des Preises von 0 CHF - 2.5 Mio CHF}
  \label{fig:verteilung_price_part}
\end{subfigure}
\caption[Verteilung des Kaufpreises]{Verteilung des Kaufpreises}
\label{fig:price}
\end{figure}

\begin{table*}[ht]
\centering
\ra{1.3}
\resizebox{\textwidth}{!}{
\begin{tabular}{@{}lrrrrr@{}}
\toprule
Objektart & Min & Max & Durchschnitt & Median & Standardabweichung\\
\midrule
Wohnung & 10 & 20'000'000 & 891’479 & 690’000 & 786’047\\
Haus & 10 & 20'000'000 & 1'281'735 & 920'000 & 1'309'085\\
Wohnung \& Haus & 10 & 20’000'000 & 1’059’503 & 790’000 & 1’066’372\\
\bottomrule
\end{tabular}}
\caption{Statistische Werte des Kaufpreises}
\label{tab:price}
\end{table*}

\subsubsection{Numerische Features}
Bis auf die Lärmbelastung wurden alle numerischen Features von den Immobilienportalen gesammelt.\\
Abbildung \ref{fig:num_features} zeigt einen Teil dieser Features in Verbindung zum Kaufpreis. Dabei ist ersichtlich, dass die Wohnfläche, Kubatur und Anzahl Zimmer eine linear ähnliche Abhängigkeit besitzen. Die restlichen Features beistzen keine direkte Korrelation zum Preis.

\begin{figure}[h]
\centering
\includegraphics[width=\textwidth]{images/Vergleich_zum_preis.png}
\caption[Numerische Features im Vergleich zum Preis]{Numerische Features im Vergleich zum Preis}%
\label{fig:num_features}
\end{figure}

Tabelle \ref{tab:num_features} zeigt die statistischen Werte für die numerischen Features. Die gesammelten Kennwerte sehen plausibel aus und decken sich zum Teil mit den Daten vom Bundesamt für Statistik\footnote{https://www.bfs.admin.ch/bfs/de/home/statistiken/bau-wohnungswesen.html}. Auffallend ist, dass der Durchschnitt und der Median bei vielen Features sehr nahe beieinander liegen. Somit gleicht die Verteilung einer Normalverteilung.

\begin{table*}[h]
\centering
\ra{1.3}
\resizebox{\textwidth}{!}{
\begin{tabular}{@{}lrrrrr@{}}
\toprule
 & Min & Max & Durchschnitt & Median & Standardabweichung\\
\midrule
Anz. Zimmer & 0& 20 & 4.84 & 4.5 & 1.9\\
Fläche & 0 & 4’000 & 143.08 & 130 & 101.924\\
Baujahr & 1300 & 2019 & 1988 & 2006 & 51.52\\
Anz. Stockwerke & 0 & 18 & 1.8 & 2 & 1.747\\
Kubatur & 1 & 18649 & 1059.74 & 867 & 838.419\\
Stockwerk & -1 & 23 & 1.25 & 1 & 1.57\\
Lärmbelastung & 0 & 73.99 & 48.8 & 49.35 & 6.85\\
Letzte Renovation & 1780 & 2019 & 2009 & 2012 & 9.09\\
\bottomrule
\end{tabular}}
\caption{Statistische Werte der numerischen Features}
\label{tab:num_features}
\end{table*}

Abbildung \ref{fig:heatmap} zeigt eine Heatmap der einzelnen Features. Diese zeigt auf, welche Features  untereinander korrelieren. Der Kaufpreis hat die höchste Korrelationen mit der Wohnfläche und der Anzahl Zimmer. Weniger starke Korrelationen besitzt er mit dem Baujahr sowie Stockwerk.

\begin{figure}[h!]
\centering
\includegraphics[width=0.9\textwidth]{images/Heatmap_numerical.png}
\caption[Heatmap der numerischen Features]{Heatmap der numerischen Features}
\label{fig:heatmap}
\end{figure}

% TODO newpage
\newpage
\subsubsection{Kategorische Features}
Bei den kategorischen Features handelte es sich vor allem um ortsbezogene Features vom BFS. Viele davon beschrieben die Region, indem sich die Immobilie befand.\\
Insgesamt wurden 35 kategorische Features verwendet. Eine Analyse der ortsbezogenen Features mit dem Kaufpreis ergab keine unerwarteten Ergebnisse. So zeigt Abbildung \ref{fig:cantons}, dass der Kanton Genf die teuersten Immobilien in der Schweiz besitzt, gefolgt vom Kanton Zug. Auch ist ersichtlich, dass städtische Immobilien im Schnitt teurer sind als auf dem Land. Anders gesagt, je urbaner die Region, desto teurer der Immobilienpreis. Diverse weitere Vergleiche befinden sich im Anhang.\\[2ex]
\begin{figure}[ht]
\centering
\includegraphics[width=\textwidth]{images/boxPlot_cantons.png}
\caption[Kaufpreis auf Kantonte aufgeteilt]{Kaufpreis auf Kantonte aufgeteilt}%
\label{fig:cantons}
\end{figure}

\textbf{Beschreibung / Merkmale}\\
Neben den ortsbezogenen Daten, analysierten wir die Beschreibung und die Merkmale der Inserate, um verschiedene Tags zuordnen zu können. Die 50 häufigsten vorkommenden Tags die in unserem Datensatz vorkamen, werden in Abbildung \ref{fig:tags} gezeigt. Die verwendeten Tags beschrieben vor allem die Innen- und Ausseneinrichtung einer Immobile.
%
\begin{figure}[ht]
\centering
\includegraphics[width=\textwidth]{images/tags.png}
\caption[Die 50 häufigsten verwendeten Tags]{Die 50 häufigsten verwendeten Tags}%
\label{fig:tags}
\end{figure}

\subsection{Machine Learning}
Insgesamt haben wir neun verschiedene Machine Learning Algorithmen untersucht. Um eine vergleichnare Bewertung zu erhalten, bekamen alle Algorithmen die gleichen Datensätze zum Trainieren und zum Testen. So wurde verhindert, dass ein Algorithmus anhand eines vorteilhafteren Datensatzes eine bessere Performance erzielte. 

Als erstes bereiteten wir unsere gesammelten Daten für unsere Algorithmen vor. Dafür wurden die Features, wenn nötig, in nummerische Werte transformiert. Dies geschah im Feature Engineering. In den nächsten Absätzen werden die verschiedenen Algorithmen mit einem Basisdatensatz untersucht.

Als Basis für unseren Datensatz transformierten wir unsere kategorischen Features mit Hilfe einer One-Hot-Encoding in binäre Features um. So enstanden aus 35 Features über 2000 Features.\\
Wie im vorherigen Kapitel gezeigt, sind die Features Kubatur, Raumhöhe, effektive Fläche, Grundstücksfläche, Anzahl Stockwerke, Stockwerk und Renovationsjahr nur selten angegeben. Somit konnten sie nicht verwendet werden. Dafür konnten die Wohnfläche, die Anzahl Zimmer, das Baujahr, das Renovationsjahr, die Beschreibung und die Merkmale einer Immobile verwendet werden.\\
Als letzter Schritt wurden die Inserate gelöscht, die keinen Eintrag bei einem verwendeten Feature besassen oder doppelt vorkamen.

\subsubsection{Regressions Algorithmen}
Für die Auswertung der Algorithmen starteten wir mit den Regressions Algorithmen. Hierfür verwendeten wir eine  einfache lineare Regression, den Ridge- und den Lasso-Algorithmus. Die Ergebnisse werden in Tabelle \ref{tab:regression} aufgezeigt.

\begin{table*}[h]
\centering
\ra{1.3}
\resizebox{\textwidth}{!}{
\begin{tabular}{@{}lrrrrr@{}}
\toprule
ML Algorithmus & $R^2$ & MAPE & MdAPE & 10\% Abweichung & Maximaler Fehler\\
\midrule
Lineare Regression & -4.26E+25 & 2.93E+13 & 21.79 & 25.03 & 5.81E+20\\
Ridge Regression & 0.33 & 61.72 & 28.82 & 19.00 & 2.87E+07\\
Lasso Regression & 0.58 & 70.04 & 19.46 & 27.82 & 2.71E+07\\
\bottomrule
\end{tabular}}
\caption{Ergebnisse der Regression Algorithmen}
\label{tab:regression}
\end{table*}

Wir sehen, dass die Algorithmen eher schlecht abschnitten. Zum einen liegt das daran, dass noch kein richtiges Feature Engineering durchgeführt wurde. Zum anderen haben wir in unserem Modell sehr viele kategorischen Features. Die lineare Regression erwartet eine Normalverteilung der Features, was bei kategorischen Features nicht der Fall ist.\\
Es ist ersichtlich, dass die Algorithmen mit einem Regularisierungsparameter eine deutlich bessere Performance erzielten, als die einfache lineare Regression. Somit ist ein Regularisierungsparameter bei diesem Datenset nützlich.\\
Auffallend war, dass der Lasso Algorithmus eine deutliche bessere Performance als der Ridge Algorithmus vorwies. Denn der Lasso kann unnötige Features eliminieren und somit die Komplexität reduzieren, während der Ridge immer alle Features verwenden muss. Abbildung \ref{fig:lasso} zeigt den Tukey-Anscombe Plot \ref{fig:lasso_tukey-anscombe} und die Verteilung der Residuen \ref{fig:lasso_histo}. Darauf ist ersichtlich, dass es extreme Ausreisser gibt und somit die MAPE verzerren.

\begin{figure}[h]
\begin{subfigure}{.5\textwidth}
  \centering
  \includegraphics[width=\linewidth]{images/lasso_tukey.png}
  \caption[Tukey-Anscombe]{Tukey-Anscombe}
  \label{fig:lasso_tukey-anscombe}
\end{subfigure}
\begin{subfigure}{.5\textwidth}
  \centering
  \includegraphics[width=\linewidth]{images/lasso_residuen.png}
  \caption[Histogramm]{Histogramm} % spaces needed for correct positioning
  \label{fig:lasso_histo}
\end{subfigure}
\caption[Auswertung der Residuen für Lasso]{Auswertung der Residuen für Lasso}
\label{fig:lasso}
\end{figure}

\subsubsection{$K$-Nearest Neighbour}
\label{chapter:KNN}
Beim nächsten Algorithmus handelte es sich um den $K$-Nearest Neighbour. Dieser Algorithmus, versucht anhand ähnlicher Inserate den Preis zu schätzen. Der Vorteil gegenüber den Regressions Algorithmen ist, dass es dem Algorithmus egal ist, in welchem Format die Features vorliegen. Zudem haben viele Features keinen negativen Einfluss auf den Algorithmus.
Tabelle \ref{tab:regression_knn} zeigt die erzielten Resultate an. Für die Berechnung wurden zwei Nachbarn zum Vergleich, eine Leaf size von 100 und für die Gewichtung die Euklid-Distanz verwendet.

\begin{table*}[ht]
\centering
\ra{1.3}
\resizebox{\textwidth}{!}{
\begin{tabular}{@{}lrrrrr@{}}
\toprule
ML Algorithmus & $R^2$ & MAPE & MdAPE & 10\% Abweichung & Maximaler Fehler\\
\midrule
$K$-Nearest Neighbour & 0.83 & 23.67 & 0 & 79.94 & 1.81E+07\\
\bottomrule
\end{tabular}}
\caption{Ergebnisse von $K$-Nearest Neighbour}
\label{tab:regression_knn}
\end{table*}

Die Performance wurde deutlich besser als bei den Regressions Algorithmen. 79,94\% der Immobilien konnten mit einer Abweichung von 10\% richtig geschätzt werden. Anhand diesen Ergebnissen konnte aufgezeigt werden, dass bei ähnlichen Immobilien auch ein ähnlicher Preis existiert.\\
Eine MdAPE von 0\% erklärten wir damit, dass es sehr ähnliche Häuser gibt, die den selben Preis besitzen. Wir gehen im Kapitel 5 nochmals genauer darauf ein. Die MAPE mit einem Wert von 23,67\% war zu hoch, um von einem guten Modell zu sprechen. Abbildung \ref{fig:kneighbour} zeigt die Statistik für die Residuen dar. Im Histogramm \ref{fig:kneighbour_histo} wird nochmals der MdAPE von 0\% deutlich. Auf dem Tukey-Anscombe Plot \ref{fig:kneighbour_tukey-anscombe} wichen die Residuen weniger stark vom beobachteten Preis ab als beim Lasso.

\begin{figure}[h]
\begin{subfigure}{.5\textwidth}
  \centering
  \includegraphics[width=\linewidth]{images/KNeighbour_tukey.png}
  \caption[Tukey-Anscombe]{Tukey-Anscombe}
  \label{fig:kneighbour_tukey-anscombe}
\end{subfigure}
\begin{subfigure}{.5\textwidth}
  \centering
  \includegraphics[width=\linewidth]{images/KNeighbour_residuen.png}
  \caption[Histogramm]{Histogramm}
  \label{fig:kneighbour_histo}
\end{subfigure}
\caption[Auswertung der Residuen für KNN]{Auswertung der Residuen für KNN}
\label{fig:kneighbour}
\end{figure}

\subsubsection{Baum Algorithmen}
Als nächstes untersuchten wir verschiedene Baum Algorithmen. Baum Algorithmen erfreuen sich aktuell einer grosser Beliebtheit bei diversen Machine Learning Wettbewerben\footnote{https://www.kaggle.com/c/house-prices-advanced-regression-techniques}.\\[2ex]
%
Wir haben vier verschiedene Algorithmen untersucht. Sie unterschieden sich hauptsächlich in ihrer Strategie. Der Random Forest wendet eine Art Bagging an. Der XGBoost wie auch der AdaBoost verwenden eine Boostingstrategie und der Extra Trees verwendet weder Boosting noch Bagging.

Eine Auswertung mittels GridSearchCV hatte für die Algorithmen folgende Parameter vorgeschlagen.
Für Random Forest verwendeten wir 700 Schätzungsmodelle, ansonsten die Default Parameter von Scikit Learn. Dasselbe galt auch für den Extra Trees Algorithmus.\\
Für AdaBoost haben wir als Basis Schätzungsmodell den DecisionTreeRegressor mit 18 Schätzungsmodelle benutzt.\\
Der XGBoost besitzte eine maximale Tiefe von 100 mit einer Lernrate von 0.1 und 350 Schätzungsmodelle.

\begin{table*}[ht]
\centering
\ra{1.3}
\resizebox{\textwidth}{!}{
\begin{tabular}{@{}lrrrrr@{}}
\toprule
ML Algorithmus & $R^2$ & MAPE & MdAPE & 10\% Abweichung & Maximaler Fehler\\
\midrule
Random Forest & 0.907 & 23.31 & 2.43 & 77.52 & 1.26E+07\\
Extra Trees & 0.91 & 23.88 & 0.86 & 81.74 & 1.04E+07\\
XGBoost & 0.898 & 18.71 & 0.81 & 82.01 & 1.46E+07\\
AdaBoost & 0.907 & 18.14 & 0.00 & 83.69 & 1.63E+07\\
\bottomrule
\end{tabular}}
\caption{Ergebnisse der Baum Algorithmen}
\label{tab:first_round}
\end{table*}

In Tabelle \ref{tab:first_round} ist erkennbar, dass die Baum Algorithmen, bis auf den Random Forest, eine noch bessere Performance hatten, als der $K$-Nearest Neighbour. Die vielen kategorischen Features halfen den Baumalgorithmen eine gute Performance zu erzielen. Wie beim $K$-Nearest Neighbour war die MAPE zu hoch, um von einem guten Modell zu sprechen.\\
Abbildung \ref{fig:adaboost_1} zeigt die Resultate für AdaBoost, da dieser die besten Resultate erzielte. Gros­so mo­do waren die Resultate ähnlich dem KNN. Wie im Tukey-Anscombe Plot \ref{fig:ada_tukey-anscombe_1} dargestellt, existierten weniger starke Ausreisser und die Residuen waren näher am Beobachtungswert. Das Histogramm \ref{fig:ada_histo_1} unterschied sich nicht gross vom KNN.

\begin{figure}[h]
\begin{subfigure}{.5\textwidth}
  \centering
  \includegraphics[width=\linewidth]{images/adaboost_tukey_anscombe_1.png}
  \caption[Tukey-Anscombe]{Tukey-Anscombe}
  \label{fig:ada_tukey-anscombe_1}
\end{subfigure}
\begin{subfigure}{.5\textwidth}
  \centering
  \includegraphics[width=\linewidth]{images/adaboost_verteilung_residuals_log_1.png}
  \caption[Histogramm]{Histogramm}
  \label{fig:ada_histo_1}
\end{subfigure}
\caption[Auswertung der Residuen für AdaBoost]{Auswertung der Residuen für AdaBoost}
\label{fig:adaboost_1}
\end{figure}

Für das weitere Vorgehen wurden nur noch der Extra Trees, XGBoost und der AdaBoost verwendet, da diese die besten Werte erreichten.
% - - - - - - - - - - - - - -
% Feature engineering
% - - - - - - - - - - - - - -
\subsubsection{Feature Engineering}
Um eine bessere Performance zu erhalten, wurde anhand von Feature Engineering versucht weitere Features zu generieren. Aus den beiden Features Wohnfläche und Anzahl Zimmer, berechneten wir die durchschnittliche Zimmergrösse. Wir nahemn diese beiden Features, da diese eine hohe Korrelation mit dem Preis besassen. Tabelle \ref{tab:second_round} zeigt die Ergebnisse auf.

\begin{table*}[ht]
\centering
\ra{1.3}
\resizebox{\textwidth}{!}{
\begin{tabular}{@{}lrrrrr@{}}
\toprule
ML Algorithmus & $R^2$ & MAPE & MdAPE & 10\% Abweichung & Maximaler Fehler\\
\midrule
Extra Trees & 0.903 & 27.14 & 0.82 & 81.426 & 1.76E+07\\
XGBoost & 0.888 & 25.27 & 0.752 & 81.94 & 2.30E+07\\
AdaBoost & 0.885 & 23.28 & 0 & 83.97 & 2.33E+07\\
\bottomrule
\end{tabular}}
\caption{Ergebnisse mit durchschnittlicher Zimmergrösse}
\label{tab:second_round}
\end{table*}

Erstaunlicherweise konnte sich kein Algorithmus verbessern. Im Gegenteil, die Schätzungen wurden schlechter, obwohl die durchschnittliche Zimmergrösse als drittwichtigstes Feature gewichtet wurde. Die maximale Abweichung von 10\% ist jedoch gleich gut geblieben. Aus diesem Grund verwendeten wir dieses Feature weiterhin für die Erstellung unsers Modells. Die höhere MAPE erklärten wir uns durch den höheren maximalen Fehler.\\
Interessanterweise besassen Wohnfläche, Anzahl Zimmer und durchschnittliche Zimmergrösse eine hohe Gewichtung, obwohl sie untereinander eine starke Korrelation aufwiesen. Normalerweise erhält nur ein Feature von mehreren stark korrelierende Features eine hohe Gewichtung.\\[2ex]
%
\textbf{Renovationsjahr}\\
Als nächstes ergänzten wir die bestehenden Features durch ein binäres Feature, das auf 1 steht, falls eine Renovation durchgeführt wurde und andernfalls auf 0 steht.\\
Zusätzlich berechnetetn wir die die Anzahl Jahre, seit der letzten Renovation. Fehlte das Renovationsjahr, wurde stattdessen das Baujahr verwendet. Durch die Transformation fiel das Renovationsjahr weg. Die Ergebnisse sind in Tabelle \ref{tab:third_round} dargestellt.

\begin{table*}[ht]
\centering
\ra{1.3}
\resizebox{\textwidth}{!}{
\begin{tabular}{@{}lrrrrr@{}}
\toprule
ML Algorithmus & $R^2$ & MAPE & MdAPE & 10\% Abweichung & Maximaler Fehler\\
\midrule
Extra Trees & 0.915152 & 20.048 & 0.87 & 81.804 & 8.87E+06\\
XGBoost & 0.910506 & 20.191 & 0.802 & 82.113 & 9.19E+06\\
AdaBoost & 0.920327 & 18.484 & 0 & 83.528 & 9.75E+06\\
\bottomrule
\end{tabular}}
\caption{Ergebnisse mit Einbezug der Renovation}
\label{tab:third_round}
\end{table*}

Obwohl das Renovationsjahr keine starke Korrelation mit dem Kaufpreis besass, konnten diese beiden Features die Performance unserer Modelle verbessern. Zusätzlich konnte der maximale Fehler um einiges verringert werden, was einen positiven Einfluss auf die MAPE hatte. Die maximale Abweichung von 10\% ist jedoch gleich geblieben.

\textbf{Lärmbelastung}\\
Für die Berechnung der Lärmbelastung nahmen wir den Durchschnitt von 255 Pixel, was etwa einem Radius von 10m entspricht. Fehlte die exakte Wohnadresse, verwendeten wir die Koordinaten der Gemeinde.\\
Abbildung \ref{fig:noise} zeigt einen Ausschnitt der Lärmkarte für eine Berechnung der Lärmbelastung. Tabelle \ref{tab:fourth_round} zeigt die Ergebnisse der Modelle, mit dem neuen Feature Lärmbelastung.

\begin{figure}[ht]
\centering
\includegraphics[width=0.5\textwidth]{images/noise.jpeg}
\caption[Ausschnitt der Lärmbelastung am Brausebad (Basel)]{Ausschnitt der Lärmbelastung am Brausebad (Basel)}
\label{fig:noise}
\end{figure}

\begin{table*}[ht]
\centering
\ra{1.3}
\resizebox{\textwidth}{!}{
\begin{tabular}{@{}lrrrrr@{}}
\toprule
ML Algorithmus & $R^2$ & MAPE & MdAPE & 10\% Abweichung & Maximaler Fehler\\
\midrule
Extra Trees & 0.838332 & 9.912 & 0.85 & 82.002 & 2.79E+07\\
XGBoost & 0.814466 & 9.952 & 0.74 & 82.212 & 2.81E+07\\
AdaBoost & 0.823295 & 9.954 & 0 & 83.593 & 2.75E+07\\
\bottomrule
\end{tabular}}
\caption{Ergebnisse mit Einbezug der Lärmbelastung}
\label{tab:fourth_round}
\end{table*}

In diesem Durchlauf hatte sich die MAPE beträchtlich verbessert, obwohl der maximale Fehler sehr hoch ist. Die MdAPE sowie die maximale Abweichung von 10\% wurden leicht besser.

\newpage
\textbf{Outlier Detection}\\
Auffallend war, dass bei allen Modellen der hohe maximale Fehler zwischen 10 und 25 Millionen lag. Um diesen Fehler zu verkleinern führten wir vor dem Lernen eine Outlier Detection durch. Somit wurden Immobilien mit abnormalen Werten aus dem Datensatz entfernt.\\
Für die Outlier Detection wendeten wir den Isolation Forest auf unsere numerischen Features an. So konnten wir fehlerhafte Inserate erkennen und entfernen.\\[2ex]
%
Beim Isolation Forest definiert ein Prozentsatz die Anzahl an markierten Ausreisser.
Um einen geeigneten Prozentsatz zu bestimmen, verglichen wir die Entwicklung der Standardabweichung bei Erhöhung des Prozentsatzes.\\
War die Differenz der Standardabweichung mit der Vorherigen nicht um mehr als 2.5\% gesunken, hatten wir unseren optimalen Prozentsatz gefunden.\\
Die 2.5\% ergaben bei uns die richtige Mischung aus genügend erkannten Ausreissern, ohne dass zu viele Inserate entfernt wurden.
Abbildung \ref{fig:outlier} zeigt die Reduktion der Standardabweichung \ref{fig:isolation_std} und die gefundenen Outlier \ref{fig:isolation_forest} bei der Wohnfläche. Tabelle \ref{tab:iso_forest} zeigt die gefunden Ausreisser im Datensatz. Insgesamt wurden 12'362 Inserate als Ausreisser markiert und entfernt. 
Es ist unschwer zu erkennen, dass gewisse Inserate nicht nur bei einem Feature als Ausreisser markiert wurden, sondern gleich bei mehreren.

\begin{figure}[ht]
\begin{subfigure}{.5\textwidth}
  \centering
  \includegraphics[width=\linewidth]{images/living_area_IsolationForest.png}
  \caption[Isolation Forest angewendet auf Wohnfläche]{Isolation Forest angewendet auf Wohnfläche}
  \label{fig:isolation_forest}
\end{subfigure}
\begin{subfigure}{.5\textwidth}
  \centering
  \includegraphics[width=\linewidth]{images/living_area_std.png}
  \caption[Entwicklung der Standardabweichung]{Entwicklung der Standardabweichung\\ \ } % spaces needed for correct positioning
  \label{fig:isolation_std}
\end{subfigure}
\caption[Outlier Detection: Wohnfläche]{Outlier Detection: Wohnfläche}
\label{fig:outlier}
\end{figure}

\begin{table*}[ht]
\centering
\ra{1.3}
\begin{tabular}{@{}lrr@{}}
\toprule
Feature & Entfernt in \% & \# entfernte Inserate\\
\midrule
Baujahr & 5 & 4293\\
Anzahl Zimmer & 3 & 2579\\
Wohnfläche & 3 & 2577\\
Letzte Konstruktion & 5 & 4294\\
Lärmbelastung & 7 & 7726\\
\bottomrule
\end{tabular}
\caption{Übersicht der entfernten Inserate mit dem Isolation Forest}
\label{tab:iso_forest}
\end{table*}
%
\begin{table*}[ht]
\centering
\ra{1.3}
\resizebox{\textwidth}{!}{
\begin{tabular}{@{}lrrrrr@{}}
\toprule
ML Algorithmus & $R^2$ & MAPE & MdAPE & 10\% Abweichung & Maximaler Fehler\\
\midrule
Extra Trees & 0.947074 & 5.216 & 0.702 & 84.269 & 1.89E+06\\
XGBoost & 0.941227 & 5.392 & 0.604 & 84.208 & 1.95E+06\\
AdaBoost &0.940529 & 4.692 & 0 & 85.609 & 1.95E+06\\
\bottomrule
\end{tabular}}
\caption{Ergebnisse mit Einbezug einer Outlier Detection}
\label{tab:fifth_round}
\end{table*}

Wird die Einführung der Outlier Detection auf die Performance der Algorithmen in Tabelle \ref{tab:fifth_round} betrachtet, ist erkennbar, dass bei allen Performancekriterien eine Verbesserung erzielt werden konnte. Der maximale Fehler wurde auf 2~Mio.~CHF verkleinert, da Extremwerte nicht mehr vorhanden waren. Deshalb ist eine Outlier Detection für ein Machine Learning Algorithmus essentiell. Dies ist auch in Abbildung \ref{fig:adaboost_6} beim Tukey-Anscombe Plot \ref{fig:ada_tukey-anscombe_6} wie auch beim Histogramm der Residuen \ref{fig:ada_histo_6} ersichtlich.

Die Tabelle \ref{tab:price_outlier} zeigt die neuen statistischen Werte für den Kaufpreis einer Immobilie. Im Vergleich zu Tabelle \ref{tab:price} ist ersichtlich, dass die Standardabweichung beträchtlich vermindert wurde. Ebenso ergeben der minimale wie auch der maximale Preis mehr Sinn. Der Durchschnitt und der Median besitzen ähnliche Werte wie vor der Outlier Detection.

\begin{figure}[h]
\begin{subfigure}{.5\textwidth}
  \centering
  \includegraphics[width=\linewidth]{images/adaboost_tukey_anscombe_6.png}
  \caption[Tukey-Anscombe]{Tukey-Anscombe}
  \label{fig:ada_tukey-anscombe_6}
\end{subfigure}
\begin{subfigure}{.5\textwidth}
  \centering
  \includegraphics[width=\linewidth]{images/adaboost_verteilung_residuals_log_6.png}
  \caption[Histogramm]{Histogramm}
  \label{fig:ada_histo_6}
\end{subfigure}
\caption[Auswertung der Residuen für AdaBoost mit Outlier Detection]{Auswertung der Residuen für AdaBoost  mit Outlier Detection}
\label{fig:adaboost_6}
\end{figure}

\begin{table*}[ht]
\centering
\ra{1.3}
\resizebox{\textwidth}{!}{
\begin{tabular}{@{}lrrrrr@{}}
\toprule
Objektart & Min & Max & Durchschnitt & Median & Standardabweichung\\
\midrule
Wohnung & 90000 & 3'500'000 & 837'394 & 705'000 & 488'630\\
Haus & 125'000 & 3'750'000 & 1'082'756 & 930'000 & 541'671\\
Wohnung \& Haus & 90000 & 3'750'000 & 931'671 & 795'000 & 523'448\\
\bottomrule
\end{tabular}}
\caption{Statistische Werte des Kaufpreises nach der Outlier Detection}
\label{tab:price_outlier}
\end{table*}

\begin{table*}[ht]
\centering
\ra{1.3}
\resizebox{\textwidth}{!}{
\begin{tabular}{@{}lrrrrr@{}}
\toprule
ML Algorithmus & $R^2$ & MAPE & MdAPE & 10\% Abweichung & Maximaler Fehler\\
\midrule
Extra Trees & 0.946107 & 5.39 & 0.686 & 83.799 & 1.55E+06\\
XGBoost & 0.940157 & 5.382 & 0.632 & 83.997 & 1.82E+06\\
AdaBoost &0.943958 & 4.64 & 0 & 85.596 & 1.81e+06\\
\bottomrule
\end{tabular}}
\caption{Ergebnisse mit Einbezug des Steuerfuss}
\label{tab:sixth_round}
\end{table*}

\newpage
\clearpage
\textbf{Steuerfuss}\\
Von Interesse war, ob der Steuersatz eine Auswirkung auf das Modell besitzt. Deshalb hatten wir vom Bundesamt für Statistik den Steuerfuss für Gemeinde und Kanton beschafft. Tablle \ref{tab:sixth_round} zeigt, dass sich weder eine Verbesserung noch eine Verschlechterung mit diesen beiden Features ergab. Der Steuerfuss besass demnach keinen grossen Einfluss auf den Immobilienpreis, beziehungsweise auf die Schätzungsmodelle.

\textbf{Tag Gruppierung}\\
Unser Modell besitzt insgesamt 50 verschiedene Tags als einzelne Feature. Diese Tags hatten sich stellenweise bedeutungsmässig überschnitten. Daher lohnte es sich aus überschneidenden Features ein Feature zu bilden. Wir definierten 4 Gruppen denen wir die Tags zuordneten. 

\begin{description}
\item[Badezimmer:] Alle Tags die eine Nasszelle beschrieben. Dazu gehörte Badezimmer, Lavabo, Dusche usw.
\item[Innenausstattung:] Alle Tags die eine Innenausstattung beschrieben, wie TV Anschluss, gross, hell usw.
\item[Aussenausstattung:] In diese Kategorie fielen Tags wie Garten, Balkon, Garage, Parkplatz usw.
\item[Nachbarschaft:] Beschrieb die Umgebung der Immobilie. Ist es zentral, ruhig, hat es Einkaufsmöglichkeiten usw.
\end{description}

\begin{table*}[ht]
\centering
\ra{1.3}
\resizebox{\textwidth}{!}{
\begin{tabular}{@{}lrrrrr@{}}
\toprule
ML Algorithmus & $R^2$ & MAPE & MdAPE & 10\% Abweichung & Maximaler Fehler\\
\midrule
Extra Trees & 0.947874 & 4.854 & 0.109 & 85.541 & 2.54E+06\\
XGBoost &0.945822 & 4.881 & 0.074 & 85.759 & 2.50E+06\\
AdaBoost &0.93973 & 4.544 & 0 & 86.317 & 2.60E+06\\
\bottomrule
\end{tabular}}
\caption{Ergebnisse mit gruppierten Tags}
\label{tab:seventh_round}
\end{table*}

Die Ergebnisse in Tabelle \ref{tab:seventh_round} zeigen nochmals eine kleine Verbesserung der Performance. In Abbildung \ref{fig:adaboost_8} ist ersichtlich, dass die geschätzten Werte nochmal ein Stück näher am Beobachtungswert liegen.

\begin{figure}[h]
\begin{subfigure}{.5\textwidth}
  \centering
  \includegraphics[width=\linewidth]{images/adaboost_tukey_anscombe_8.png}
  \caption[Tukey-Anscombe]{Tukey-Anscombe}
  \label{fig:ada_tukey-anscombe_8}
\end{subfigure}
\begin{subfigure}{.5\textwidth}
  \centering
  \includegraphics[width=\linewidth]{images/adaboost_verteilung_residuals_log_8.png}
  \caption[Histogramm]{Histogramm}
  \label{fig:ada_histo_8}
\end{subfigure}
\caption[Auswertung der Residuen für AdaBoost mit komplettem Feature Engineering]{Auswertung der Residuen für AdaBoost mit komplettem Feature Engineering}
\label{fig:adaboost_8}
\end{figure}

\begin{table*}[ht]
\centering
\ra{1.3}
\resizebox{\textwidth}{!}{
\begin{tabular}{@{}lrrrrr@{}}
\toprule
ML Algorithmus & $R^2$ & MAPE & MdAPE & 10\% Abweichung & Maximaler Fehler\\
\midrule
Extra Trees & 0.907327 & 7.092 & 0.684 & 79.063 & 1.96E+06\\
XGBoost & 0.892067 & 8.692 & 0.941 & 78.207 & 2.16E+06\\
AdaBoost & 0.893971 & 7.054 & 0 & 80.176 & 1.96E+06\\
\bottomrule
\end{tabular}}
\caption{Ergebnisse ohne ortsbezogenen Daten vom BFS}
\label{tab:height_round}
\end{table*}

\newpage
\clearpage

\textbf{Ohne ortsbezogene Daten}\\
Mit unserem Feature Engineering konnten wir aufzeigen, dass eine deutliche Verbesserung der Modelle erreicht werden konnte. Für die Forschungsfrage, \textit{Wie stark kann das Resultat verbessert werden, wenn geobasierte Daten hinzugefügt werden?}, untersuchten wir die Resultate ohne Daten vom BFS. Es ist in Tabelle \ref{tab:height_round} gut ersichtlich, dass ohne ortsbezogene Daten ein schlechteres Modell berechnet wurde.

\subsubsection{Kombinierte Algorithmen}
Als weiteren Test haben wir einen eigenen Regressor erstellt, der mehrere Machine Learnig Algorithmen kombiniert.

Dieser Algorithmus hatte im Endeffekt zwei Schätzphasen. Zunächst wurde ein Extra Trees trainiert, und für jeden Train-Datensatz wurden alle Decision Pathes für alle Bäume abgespeichert. Um nun eine Schätzung zu erhalten, wurden zuerst die Decision Pathes für den zu testenden Datensatz herausgesucht, und mit den Train-Datensätzen verglichen. Alle Train-Datensätze mit mindestens einem passendem Decision Path wurden markiert für die zweite Phase der Schätzung.\\
In der zweiten Phase wurden diese markierten Datensätze auf einem weiteren Machine Learning Algorithmus ad-hoc trainiert. Der zweite Algorithmus konnte nach belieben ausgetauscht werden. Anschliessend wurde auf diesem neuen Modell der Test-Datensatz geschätzt.\\
Theoretisch wäre es möglich gewesen, für alle möglichen Decision Pathes die Modells im Vorhinein zu berechnen, da nur eine finite Anzahl an Decision Pathes existierten. Dies war aber aus praktikablen Gründen nicht möglich, da mit vielen Datensätzen die trainierten Bäume gross wurden und entsprechend dafür zu viel RAM und Rechenzeit benötigt wurde.

Um genügend Datensätze für die zweite Phase zu erhalten, wurde der Extra Trees mit einer minimalen Leaf Size von fünf Samples konfiguriert.\\
Für den zweiten Schätzungsalgorithmus verwendeten wir den $K$-Nearest Neighbor wie oben im Kaptiel \ref{chapter:KNN} beschrieben. Wie in Tabelle \ref{tab:combined} erkennbar, waren die Ergebnisse weitaus schlechter als oben verwendete Algorithmen.

\begin{table*}[ht]
\centering
\ra{1.3}
\resizebox{\textwidth}{!}{
\begin{tabular}{@{}lrrrrr@{}}
\toprule
ML Algorithmus & $R^2$ & MAPE & MdAPE & 10\% Abweichung & Maximaler Fehler\\
\midrule
Extra Trees mit KNN & 43.636 & 13.451 & 40.034 &  0.619114 & 2.94294e+07\\
\bottomrule
\end{tabular}}
\caption{Ergebnisse eigenentwickelter Algorithmus}
\label{tab:combined}
\end{table*}

Als letztes kombinierten wir noch die drei Baumalgorithmen miteinander. Verschiedene Tests zeigten, das die Gewichtung von 70\% AdaBoost, 20\% XGBoost und 10\% Extra Trees die beste Performnace lieferte.
Das Resultat in Tabelle \ref{tab:combined_2} zeigt die Performancewerte.

\begin{table*}[ht]
\centering
\ra{1.3}
\resizebox{\textwidth}{!}{
\begin{tabular}{@{}lrrrrr@{}}
\toprule
ML Algorithmus & $R^2$ & MAPE & MdAPE & 10\% Abweichung & Maximaler Fehler\\
\midrule
Kombiniert & 0.94646 & 4.416 & 0.07 & 87.03	& 1.77E+06\\
\bottomrule
\end{tabular}}
\caption{Kombination von AdaBoost, XGBoost und Extra Trees}
\label{tab:combined_2}
\end{table*}

\subsection{Auswertung}
Es hat sich gezeigt, dass eine Kombination des AdaBoost, Extra Trees und XGBoost die besten Resultate erzielt (siehe Tabelle \ref{tab:combined_2}). Werden die Algorithmen alleine betrachtet, hatte bei unseren Daten der AdaBoost die besten Ergebnisse erreicht.\\
Von der Berechnungszeit ist AdaBoost klar der Beste Algorithmus, da dieser für ein Fit ungefähr 10 Minuten benötigte, während XGBoost und Extra Trees über eine Stunde brauchten.

Wir konnten mit unserer eigenen Implementation eines Machine Learning Algorithmus keinen Erfolg erzielen, da die Performance nicht zufriedenstellend war. Bei den linearen Regressionsalgorithmen konnte auch keine akzeptable Performance erreicht werden. Der Grund dafür liegt bei den vielen kategorischen Features. Der $K$-Nearest Neighbor ist zwar viel besser als die linearen Regressionen, aber kommt nicht an die Ergebnisse der Baumalgorithmen an.

Weiter konnte bewiesen werden, dass ortsbezogene Daten das Resultat masgeblich verbessern konnten. In unserem Fall büsste die MAPE 2.5\% ein und es konnten 7\% weniger Häuser mit einer maximalen Abweichung von 10\% richtig geschätzt werden. So konnte aufgezeigt werden, dass bei einer Schätzung eines Immobilienpreises, der Standort mit seinen Eigenschaften für das Schätzungsmodell wichtig war.\\[2ex]
%
Die Qualität der Daten von Immobilienplattformen war eher durchmischt. Es existierten qualitativ sehr gute Inserate, mit vielen Kennwerten und einer ausführlichen Beschreibung. Im Schnitt jedoch konnten neben dem Preis nur die vier Kennwerte Wohnfläche, Anzahl Zimmer, Baujahr und Renovationsjahr verwendet werden.\\
Mit einem ausgiebigen Feature Engineering und einer genügenden Menge an Inseraten, fielen die qualitativ schlechten Inserate weniger ins Gewicht, beziehungsweise wurden aussortiert.  So konnten von insgesamt 162’225 Inserate 85’845 für den Machine Learning Teil aufbereitet werden. Dies entspricht 52\% der Inserate.\\
Die Outlier Detection spielte dabei eine wichtige Rolle, da sie den maximalen Fehler stark reduzierte, indem sie 12'362 abnormale Inserate entfernte. Zudem reduzierte sie die Standardabweichung um mehr als die Hälfte.

Beim Crawlen zeigte sich, dass es sich lohnte über mehrere Proxyinstanzen zu gehen, um eine Sperrung zu verhindern. Die Inserate konnten weitgehend ohne Probleme gesammelt und aufbereitet werden.
%
\subsection{Gültigkeitsrisiken}
Zu beachten gilt, dass wir unsere Modelle auf geschätzten Werten trainiert haben. Somit sind unsere Schätzungen als Unterstützung und nicht als Endpreis zu betrachten. Zudem ist der Immobilienpreis immer gewissen Schwankungen ausgesetzt. Die Modelle schätzen somit die Immobilienpreise in der gesammelten Zeit. Für eine spätere Evaluation, müsste ein neues Datenset mit aktualisierten Immobilienpreise verwendet werden.\\
Ein Blick auf die Verteilung des Preises zeigt, dass sich der Preis primär zwischen 250'000 CHF und 1.2~Mio,~CHF befindet. Durch das können extraordinäre, sowie günstige  Immobilien nur sehr schwer geschätzt werden.\\[2ex]
%
Wir haben die Kennwerte von öffentlichen Immobilienplattformen gesammelt. Dort geht es hauptsächlich darum, die Immobilie zu verkaufen. So wird gerne bei denn Kennwerten geschummelt, oder schöngeredet indem sie falsch eingetragen werden. Eine Outlier Detection kann nur grobe Ausreisser erkennen, kleinere Ausreisser bleiben somit unentdeckt und verfälschen das Modell.\\
Wird eine neue Siedlung oder ein neuer Wohnblock inseriert, kann es vorkommen, dass 10 mal die gleiche Wohnung inseriert wird. Dies ist unschön, da diese Wohnungen im Modell mehr Gewicht erhalten. Wir haben uns jedoch stark bemüht, solche Duplikate zu erkennen und zu entfernen.

%1970 von ingesamt 73483
\newpage
\clearpage
%
\section{Diskussion}
In dieser Arbeit konnte erfolgreich aufgezeigt werden, dass anhand von öffentlichen Daten und modernen Machine Learning Methoden eine akkurate Schätzung abgeben werden kann.\\[2ex]
%
Das sammeln von öffentlichen Daten ist jedoch mit einem gewissen Aufwand verbunden. Die Datenqualität der verschiedenen Immobilienplattformen sind eher dürftig. Es konnten somit nur sehr wenige Kennwerte einer Immobilie verwendet werden. Bei einer professionellen Schätzung kämen noch diverse andere Kennwerte zum Einsatz.\\
Ist die Idee über längere Zeit Daten zu sammeln, lohnt es sich in einen autonomen Crawler Zeit zu investieren. So ist man nicht abhängig vom XPath oder CSS und ist gegenüber einem Redesign der Seiten gewappnet. Eine weitere Möglichkeit wäre die Daten bei einem vorhandenen Immobilienplattform direkt zu beziehen.\\
Die Verknüpfung der ortsbezogenen Daten mit den Inseraten kann zum Teil mühsam sein. Wir haben OpenStreetMap verwendet, um möglichst Kostengünstig die Koordinaten zu erhalten. Der Nachteil an OpenStreetMap ist, dass es mehrheitlich nur Städte abdeckt. Hier lohnt sich auf jeden Fall ein Konto bei Google Maps, da die Adresssuche viel besser ist.\\[2ex]
%
Das Crawlen über ein Proxy lohnt sich, da somit die Gefahr, als Crawler erkannt zu werden, verringert wird. Durch das wir über einen Server ausserhalb der Schweiz gingen, konnten wir nicht alle Plattformen, wie zum Beispiel Comparis, aufrufen.\\
Es ist eventuell sogar sinnvoll nur von einer Plattform die Daten zu sammeln. So müssen nicht die Unterschiedlichen Bezeichnungen harmonisiert werden.\\[2ex]
%
Das Untersuchen eines Machine Learning Algorithmus ist nicht immer einfach. Vor allem bei Baumalgorithmen. Denn die Teilung ist nicht immer gleich, auch wenn immer der gleiche Seed verwendet wird. So kann es sein, dass bei einem erneuten Durchlauf ein wenig bessere oder schlechtere Resultate herauskommen.\\
Zusätzlich muss beachtet werden, dass wir vor allem Standardimmobilien in unserem Datensatz haben. Sprich ausgefallene Inserate können nur schwer geschätzt werden, wenn sie nicht schon bei der Outlier Detection herausgenommen wurden. Somit eignen sich unsere Modell nicht für ausgefallene Immobilien.\\[2ex]
%
Aufgefallen ist uns auch, dass je mehr Inserate man verwendet desto besser wurden die Baumalgorithmen. Verwendeten wir nur die Hälfte der Inserate konnten die Modelle nur noch etwa 72\% Prozent mit einer Abweichung von 10\% richtig schätzen. Besitzt man weniger Inserate, könnte ein Baumalgorithmus nicht unbedingt am geeignesten sein. Hier lohnt es sich eventuell den K-Nearest Neighbour zu verwenden.
\\[2ex]
%
Unsere eigen aufgebautes Pipeline-System hat uns beim Iterativen Prozess sehr unterstützt. So konnte ohne grosse Mühe diverse Szenarien durchgerechnet werden, ohne dass dazwischen Daten verändert werden müssen.
%	
\subsection{Weiterführende Anwendung} 
Es gibt noch viele weitere Machine Learning Algorithmen, die man untersuchen könnte. Unter Anderem wäre der LightGBM wie auch Neurale Netzwerke interessant zu untersuchen.\\ 
Weiter können unsere Modelle für Immobilienmakler verwendet werden. Dazu müsste man eine Web-Applikation zur Verfügung stellen, bei dem die Kennwerte eingetragen werden können und der Preis geschätzt wird. 


\newpage
\clearpage
%
\section{Zusammenfassung}
% zusammenfassung Kapitel 3, 4, 5
Mittels einer eigens entwickelter Controller App werden die beiden Ansätze für regelbasierte und situationsbasierte Systeme für den Benutzertest zur Verfügung gestellt. Die App integriert sich ins Samsung SmartThings Ökosystem. Mittels einer Rule Engine werden die Regeln und Situationen ausgewertet und sofern nötig Aktionen auf IoT Geräten ausgelöst.\\[2ex]
%
Die Testpersonen wurden sukzessive an jeweils eine von drei vordefinierten Testaufgaben herangeführt und mussten diese lösen. Die Testaufgaben waren aufgeteilt nach einfacher, mittlerer und hoher Komplexität. Alle Aufgaben konnten von allen Testpersonen gelöst werden. Als Merkmal hat sich dabei hervorgehoben, dass der regelbasierte Ansatz bei einfacher Komplexität bevorzugt wurde, während bei der Aufgabe mit hoher Komplexität klar der situationsbasierte Ansatz favorisiert wurde. Die Ergebnisse werden durch Feedback der Testpersonen, sowie durch die statistische Analyse gestärkt.\\[2ex]
%
Der Benutzer sollte daher über das Einsatzgebiet und die Aufgaben des System im Klaren sein, bevor er sich ein IoT System zulegt. Für Lieferanten lohnt es sich, sich in das Thema einzulesen und zumindest experimentell neue Ansätze auszuprobieren. Mit dem Momentum das IoT zurzeit besitzt, eröffnen sich im Markt immer wieder neue Möglichkeiten für die Konfiguration komplexer Systeme.
\newpage
\clearpage
\addcontentsline{toc}{section}{Literatur}
% Literaturliste anzeigen
\bibliographystyle{ieeetr}
\bibliography{literatur}
%
\clearpage
\addcontentsline{toc}{section}{Abbildungen}
\listoffigures
%
\addcontentsline{toc}{section}{Tabellen}
\listoftables
%
\renewcommand{\lstlistlistingname}{Listingverzeichnis}
\newpage
\section{Anhang}
%
\end{document}
